% part: intuitionistic-logic
% chapter: introduction
% section: axiomatic-derivations

\documentclass[../../../include/open-logic-section]{subfiles}

\begin{document}

\olfileid[zh]{int}{int}{axd}

\olsection{公理化的\usetoken{P}{derivation}}

直觉主义命题逻辑的公理化!!{derivation}在概念上是最简单的，在历史上也是最早的!!{derivation}系统。它们的运作方式与经典命题逻辑完全一样。

\begin{defn}[!!^{derivability}]
若 $\Gamma$ 是 $\Lang L$ 的!!{formula}的一个集合，则从 $\Gamma$ 出发的一条
\emph{!!{derivation}} 是一个有穷序列 $!A_1$、
\dots,~$!A_n$（其中各项均为 $\Lang L$ 的!!{formula}），并且对每个 $i \le n$ 下列条件之一成立：
\begin{enumerate}
\item $!A_i \in \Gamma$；或
\item $!A_i$ 是一条公理；或
\item $!A_i$ 由某些 $!A_j$ 与 $!A_k$（其中 $j < i$，$k <
  i$）通过分离规则推出，即 $!A_k \ident !A_j \lif !A_i$。
\end{enumerate}
\end{defn}

\begin{defn}[公理]
直觉主义命题逻辑的 \emph{公理}的集合 $\PAx$，是所有下列形式的全部!!{formula}：
\begin{align}
  & (!A \land !B) \lif !A \ollabel{ax:land1}\\
  & (!A \land !B) \lif !B \ollabel{ax:land2}\\
  & !A \lif (!B \lif (!A \land !B)) \ollabel{ax:land3}\\
  & !A \lif (!A \lor !B) \ollabel{ax:lor1}\\
  & !A \lif (!B \lor !A) \ollabel{ax:lor2}\\
  & (!A \lif !C) \lif ((!B \lif !C) \lif ((!A \lor !B) \lif !C)) \ollabel{ax:lor3}\\
  & !A \lif (!B \lif !A) \ollabel{ax:lif1}\\
  & (!A \lif (!B \lif !C)) \lif ((!A \lif !B) \lif (!A \lif !C)) \ollabel{ax:lif2}\\
  & \lfalse \lif !A \ollabel{ax:lfalse1}
\end{align}
\end{defn}

\begin{defn}[!!^{derivability}]
把一条!!^a{formula}~$!A$ 称为从 $\Gamma$ \emph{!!{derivable}}，记为
$\Gamma \Proves !A$，当且仅当存在一条从 $\Gamma$ 出发、以
$!A$ 结尾的!!a{derivation}。
\end{defn}

\begin{defn}[定理]
若存在一条从空集出发、得出~$!A$ 的!!a{derivation}，则把一条!!^a{formula}~$!A$ 称为一条 \emph{定理}。若 $!A$ 是定理，我们记 $\Proves !A$；若它不是定理，则记 $\Proves/ !A$。
\end{defn}

\begin{prop}
  若在直觉主义逻辑的公理化!!{derivation}系统中有 $\Gamma \Proves !A$，则在
  经典逻辑中也有 $\Gamma \Proves !A$。特别是，若 $!A$ 是直觉主义
  定理，则它也是经典定理。
\end{prop}

\begin{proof}
  每条直觉主义公理同时也是经典公理，因此直觉主义逻辑中的每条!!{derivation}也是经典逻辑中的一条!!a{derivation}。
\end{proof}

\end{document}
